% !TeX root = ../../Book1.tex
% 纲要标题：## 第12章　多项式环、不可约性与有限性
% 纲要位置：计划纲要.md，第 442 行。

\chapter{多项式环、不可约性与有限性}
\label{b1:ch:12}
\BookChapterNumber{b1:ch:12}

\begin{chapterabstract}
多项式的分解把方程的根、整环的算术和理想的有限生成性联系起来。本章先用带余除法建立根与重数的理论，再以 Gauss 引理把分式域上的分解降回系数整环，证明唯一分解性在添加不定元后保持。Eisenstein 判据与模素数约化法提供具体的不可约性检验。对多元多项式，先证明对称多项式可唯一用初等对称多项式表示，再从理想升链条件出发证明 Hilbert 基定理，分别研究不变环的生成与理想的有限生成性。
\end{chapterabstract}

\begin{learninggoals}
\begin{itemize}
\item 能区分多项式、它定义的函数及其在不同域中的根，并用形式导数判别重根。
\item 掌握本原多项式与容度的作用，说明分式域上的分解为何能够下降。
\item 会选择不可约性判据，说明判据的假设与检验失败时能够得出的结论。
\item 会用初等对称多项式表示对称多项式，理解消项过程的终止性和表达式的唯一性，并能用根与系数关系作消元计算。
\item 理解 Noether 性与理想升链条件的等价性，掌握 Hilbert 基定理的消项证明。
\end{itemize}
\end{learninggoals}

本章使用\cref{b1:prop:field-poly-division}的域上多项式带余除法、\cref{b1:prop:pid-bezout}的 Bézout 等式，以及\cref{b1:prop:ufd-prime}的唯一分解整环中不可约元的素性。分式域的基本构造在需要时就地给出，不以尚未建立的局部化理论为前提。第一节始终在域上讨论，第二、三节的系数整环须满足所列的唯一分解假设；最后一节证明 Hilbert 基定理时，只要求底环交换且为 Noether 环，不要求它是整环。

\ProseParagraphBreak
若以有限域和后续 Galois 理论为目标，可先连续读\secref{b1:sec:1201}至\secref{b1:sec:1203}，掌握一元多项式的根、分解与不可约性判据；遇到多元例子时，再补读\subsecref{b1:subsec:1204:01}的记号和次数性质。\subsecref{b1:subsec:1204:invariants}与\subsecref{b1:subsec:1204:symmetric-polynomials}把根的排列同系数联系起来，可供后续一般方程的讨论使用。\subsecref{b1:subsec:1204:02}与\subsecref{b1:subsec:1204:03}研究理想的有限生成性，为后续交换代数作准备，并非前述不可约性判据的前提。章末习题中的一般系数环判据只需把多项式看作 \(k[y][x]\) 中的元素；讨论 Noether 性的题目则要用到最后两个小节。

\ProseParagraphBreak
多项式与根的基础可读柯斯特利金的《代数学引论》第一卷：基础代数（张英伯译）\cite[第5--6章]{Kostrikin2006AlgebraIChinese}，或席南华的《基础代数》第一卷\cite[第6--7章]{Xi2016BasicAlgebraI}。读完 Hilbert 基定理后，若想研究多元多项式的理想成员判定与消元，可依次读 Cox、Little 与 O’Shea 的《Ideals, Varieties, and Algorithms》\cite[Chapter 2, \S\S 2--8; Chapter 3, \S\S 1--2]{CoxLittleOShea2015IdealsVarietiesAlgorithms}：先在第2章第2--4节补上单项式序、多元除法和单项式理想，再读第5--8节的 Gröbner 基与算法，最后看第3章第1--2节的消去和几何应用。若要追踪延拓定理与闭包定理的完整证明，还须续读第3章第5--6节和第4章第7节\cite[Chapter 3, \S\S 5--6; Chapter 4, \S 7]{CoxLittleOShea2015IdealsVarietiesAlgorithms}；第3章开头的定理与应用不能代替这些证明。

% !TeX root = ../../Book1.tex
% 纲要标题：### 12.1 一元多项式
% 纲要位置：计划纲要.md，第 444 行。

\section{一元多项式}
\label{b1:sec:1201}

% TODO: 按以下纲要要点撰写本节。

% 纲要内容（仅作写作计划，尚非教材正文）：
%
% * 带余除法、根与重数
% * 形式导数与重根判别
% * 插值及代数基本定理的代数部分
%

\subsection{带余除法、根与重数}
\label{b1:subsec:1201:01}

待撰写。

\subsection{形式导数与重根判别}
\label{b1:subsec:1201:02}

待撰写。

\subsection{插值及代数基本定理的代数部分}
\label{b1:subsec:1201:03}

待撰写。

% !TeX root = ../../Book1.tex
% 纲要标题：### 12.2 Gauss 引理
% 纲要位置：计划纲要.md，第 450 行。

\section{Gauss 引理}
\label{b1:sec:1202}

% 纲要内容（仅作写作计划，尚非教材正文）：
%
% * 本原多项式与容度
% * 首次完整构造分式域，再证明分式域上的分解及其下降
% * 唯一分解整环上的多项式仍为唯一分解整环
%

整数系数多项式在有理数域中分解以后，因子的分母是否真的增加了分解的可能性？答案取决于能否先把系数的公共因子剥离出来。本节始终设 \(R\) 为唯一分解整环。前一章已经证明的“不可约元必为素元”，将使约化到 \(R/(p)\) 的办法适用。

\subsection{本原多项式与容度}
\label{b1:subsec:1202:01}

\begin{definition}\label{b1:def:polynomial-content}
非零多项式 \(f=a_0+\cdots+a_nx^n\in R[x]\) 的各系数的一个最大公因子称为它的\term[rongdu]{容度}[content]，记为 \(c(f)\)。若 \(c(f)\) 是单位，就称 \(f\) 为\term[benyuan-duoxiangshi]{本原多项式}[primitive polynomial]。容度只在相伴意义下确定；当 \(R=\mathbb Z\) 时，统一选取正的最大公因子。
\end{definition}
\symidx{\(c(f)\)：非零多项式的容度}

最大公因子存在，因为在各非零系数的素因子分解中取指数的最小值即可；零系数不增加限制。将每个系数除以 \(c(f)\)，得到 \(f=c(f)f_0\)，其中 \(f_0\) 本原。例如 \(6x^2+9x+3=3(2x^2+3x+1)\)。本原并不要求某个系数为单位：\(2x+3\) 在 \(\mathbb Z[x]\) 中本原，但两个非零系数都不是单位。

\begin{lemma}[Gauss]\label{b1:lem:gauss-primitive}
两个本原多项式的乘积仍为本原多项式。因此，对非零 \(f,g\in R[x]\)，容度 \(c(fg)\) 与 \(c(f)c(g)\) 相伴。
\end{lemma}
\begin{proof}
设 \(f,g\) 本原，而 \(fg\) 不本原。则存在不可约元 \(p\)，整除 \(fg\) 的每个系数。由\cref{b1:prop:ufd-prime}，\(p\) 是素元，所以 \(R/(p)\) 是整环：两个剩余类的积为零意味着 \(p\mid ab\)，进而 \(p\mid a\) 或 \(p\mid b\)。逐系数约化得到
\[
\bar f\bar g=0\qquad\text{于 }\Bigl(R/(p)\Bigr)[x].
\]
因 \(f,g\) 本原，\(p\) 不可能分别整除它们的全部系数，故 \(\bar f,\bar g\) 都非零。\cref{b1:prop:polynomial-degree}说明整环上的非零多项式之积非零，产生矛盾。

一般地，写 \(f=c(f)f_0\)、\(g=c(g)g_0\)。刚才证明了 \(f_0g_0\) 本原，所以从系数的素因子指数取最小值得到，\(fg=c(f)c(g)f_0g_0\) 的容度与 \(c(f)c(g)\) 相伴。
\end{proof}

这个结论称为\term[gauss-yinli]{Gauss 引理}[Gauss lemma]。证明的关键在于 \(R/(p)\) 没有零因子，而非对多项式系数逐项追踪可能发生的抵消。唯一分解性保证不可约元具有这里所需的素性。

\subsection{分式域上分解的下降}
\label{b1:subsec:1202:02}

为了在更大的域中分解多项式，先完整构造整环的分式域。下面的构造不要求唯一分解性；只有随后的约分与本原多项式论证需要这一假设。\secref{b1:sec:1301}将以这里的结果为基础，研究分式域的泛性质，再由\secref{b1:sec:1302}推广到一般局部化。

\begin{proposition}\label{b1:prop:fraction-field-basic}
每个整环 \(R\) 都嵌入一个域 \(\operatorname{Frac}(R)\)，其中每个元素写成 \(a/b\)，\(a,b\in R\)、\(b\ne0\)。两个分数相等当且仅当交叉相乘相等，运算为
\[
\frac ab+\frac cd=\frac{ad+bc}{bd},\qquad
\frac ab\frac cd=\frac{ac}{bd}.
\]
这个域称为 \(R\) 的\term[fenshi-yu]{分式域}[field of fractions]。
\end{proposition}
\symidx{\(\operatorname{Frac}(R)\)：整环 \(R\) 的分式域}
\begin{proof}
在所有有序对 \((a,b)\)、\(b\ne0\) 上定义
\((a,b)\sim(c,d)\) 当且仅当 \(ad=bc\)。自反性与对称性直接成立。若另有 \(ce=df\)，则由 \(ad=bc\) 得
\(ade=bce=bdf\)。因 \(d\ne0\)，整环的消去律给出 \(ae=bf\)，所以关系具有传递性。将等价类记为 \(a/b\)。

要验证运算良定义，设 \(ab'=a'b\)、\(cd'=c'd\)。一方面
\[
(ad+bc)b'd'=a'd'bd+b'c'bd=(a'd'+b'c')bd;
\]
另一方面 \(acb'd'=a'c'bd\)。这恰好分别说明更换两个分数的代表后，所给和与积不变。分母的积非零仍由整环假设保证。

加法交换律和乘法交换律继承自 \(R\)。任意三个分数 \(a/b,c/d,e/f\) 的两种加法结合方式都给出 \((adf+bcf+bde)/(bdf)\)，两种乘法结合方式都给出 \(ace/(bdf)\)。分配律两边分别算得
\[
\frac ab\biggl(\frac cd+\frac ef\biggr)
=\frac{a(cf+de)}{bdf},\qquad
\frac{ac}{bd}+\frac{ae}{bf}
=\frac{ab(cf+de)}{b^2df};
\]
它们按交叉相乘相等。零元为 \(0/1\)，单位元为 \(1/1\)，\(a/b\) 的加法逆元为 \((-a)/b\)。分数 \(a/b\) 非零当且仅当 \(a\ne0\)，这时 \(b/a\) 为乘法逆元。因此所得结构是域。映射 \(a\mapsto a/1\) 保持和、积与单位，且 \(a/1=0/1\) 迫使 \(a=0\)，故为所需嵌入。
\end{proof}

以下记 \(F=\operatorname{Frac}(R)\)，并通过上述嵌入把 \(R[x]\) 看作 \(F[x]\) 的子环。由于 \(R\) 唯一分解，每个非零分数均可约去分子、分母的公共素因子，写成互素的 \(a/b\)。

\begin{lemma}\label{b1:lem:primitive-scaling}
每个非零 \(h\in F[x]\) 可写成 \(h=\lambda h_0\)，其中 \(\lambda\in F\) 非零，\(h_0\in R[x]\) 本原。若 \(h_0\) 本原且 \(\lambda h_0\in R[x]\)，则 \(\lambda\in R\)；若 \(\lambda h_0\) 也本原，则 \(\lambda\) 是 \(R\) 的单位。
\end{lemma}
\begin{proof}
取 \(h\) 的有限个系数分母的乘积 \(b\ne0\)，则 \(bh\in R[x]\)。把 \(bh\) 的容度 \(a\) 提出，得到 \(h=(a/b)h_0\)，其中 \(h_0\) 本原。

现在将 \(\lambda=a/b\) 约成互素形式。若 \(\lambda h_0\) 的系数都在 \(R\)，对 \(h_0\) 的每个系数 \(u\) 都有 \(au=bv\)，其中 \(v\in R\)。若 \(b\) 不是单位，选不可约因子 \(p\mid b\)。互素性给出 \(p\nmid a\)，而素性迫使 \(p\mid u\)。这对每个系数成立，与 \(h_0\) 本原矛盾。因此 \(b\) 为单位，即 \(\lambda\in R\)。若 \(\lambda h_0\) 也本原，任何非单位 \(\lambda\) 的素因子都会整除它的全部系数，仍产生矛盾，故 \(\lambda\) 为单位。
\end{proof}

\begin{theorem}[Gauss 分解定理]\label{b1:thm:gauss-descent}
设 \(f\in R[x]\) 本原且次数为正。它在 \(R[x]\) 中不可约，当且仅当它在 \(F[x]\) 中不可约。此外，若 \(f\) 本原、\(h\in R[x]\)，且 \(f\mid h\) 在 \(F[x]\) 中成立，则该整除关系已经在 \(R[x]\) 中成立。
\end{theorem}
\begin{proof}
先证最后的整除断言。\(h=0\) 时取商为零即可；否则写 \(h=fq\)，并由\cref{b1:lem:primitive-scaling}写 \(q=\lambda q_0\)，其中 \(q_0\in R[x]\) 本原。\cref{b1:lem:gauss-primitive}说明 \(fq_0\) 本原，而 \(h=\lambda fq_0\in R[x]\)，再由\cref{b1:lem:primitive-scaling}得 \(\lambda\in R\)。所以 \(q\in R[x]\)。

若 \(f=gh\) 在 \(F[x]\) 中分解，且两因子次数均为正，将它们写成 \(g=\alpha g_0\)、\(h=\beta h_0\)，其中 \(g_0,h_0\) 是 \(R[x]\) 中的本原多项式。于是
\(f=\alpha\beta g_0h_0\)。\cref{b1:lem:gauss-primitive}说明 \(g_0h_0\) 本原，且 \(f\) 也本原，所以\cref{b1:lem:primitive-scaling}给出 \(\alpha\beta\in R^\times\)。因此 \(f=(\alpha\beta g_0)h_0\) 是 \(R[x]\) 中的非平凡分解。

反过来，若本原的 \(f\) 在 \(R[x]\) 中分解为两个非单位，由容度的乘法性质，两因子均本原。整环上多项式环的单位只能是底环的单位：若 \(uv=1\)，次数相加为零，故 \(u,v\) 都是常数。因此本原的非单位因子不能为常数，两因子均有正次数，同一分解在 \(F[x]\) 中仍非平凡。
\end{proof}

所以整数系数的本原多项式在 \(\mathbb Q[x]\) 中可约，当且仅当能拆成两个正次数的整数系数多项式。这里不能省略“本原”：常数 \(2\) 在 \(\mathbb Z[x]\) 中是非单位，在 \(\mathbb Q[x]\) 中却是单位；\(2x+2\) 的容度也必须先提出，再讨论真正涉及次数的分解。

\subsection{多项式环的唯一分解性}
\label{b1:subsec:1202:03}

\begin{theorem}\label{b1:thm:polynomial-ufd}
若 \(R\) 是唯一分解整环，则 \(R[x]\) 也是唯一分解整环。因而，对任意有限个不定元，\(R[x_1,\ldots,x_n]\) 仍为唯一分解整环。
\end{theorem}
\begin{proof}
仍令 \(F=\operatorname{Frac}(R)\)。由\cref{b1:prop:field-poly-division}、\cref{b1:thm:euclidean-pid}及\cref{b1:thm:pid-ufd}，\(F[x]\) 是唯一分解整环。对 \(R[x]\) 中非零的 \(f\)，先写 \(f=c(f)f_0\)，其中 \(f_0\) 本原。容度 \(c(f)\) 在 \(R\) 中有不可约分解。若 \(\deg f_0>0\)，再在 \(F[x]\) 中分解 \(f_0\)，把每个正次数不可约因子乘以适当非零常数，化为 \(R[x]\) 中的本原因子 \(q_1,\ldots,q_s\)。于是
\[
f_0=\lambda q_1\cdots q_s.
\]
右边的多项式乘积本原，故\cref{b1:lem:primitive-scaling}说明 \(\lambda\in R^\times\)。而\cref{b1:thm:gauss-descent}说明每个 \(q_i\) 在 \(R[x]\) 中不可约。若 \(f_0\) 为常数，则本原性已说明它为单位，不需这一步。

还要证明这些分解唯一。\(R\) 的每个不可约元 \(p\)，在 \(R[x]\) 中都是素元：模 \(p\) 约化后的多项式环 \(\Bigl(R/(p)\Bigr)[x]\) 是整环，故 \(p\mid gh\) 蕴含 \(p\mid g\) 或 \(p\mid h\)。正次数的本原因子 \(q_i\) 在 \(F[x]\) 中不可约，\cref{b1:prop:ufd-prime}说明它在这个唯一分解整环中是素元。若 \(q_i\mid gh\) 于 \(R[x]\)，先在 \(F[x]\) 中得到 \(q_i\mid g\) 或 \(q_i\mid h\)，再由\cref{b1:thm:gauss-descent}的整除下降，得到相应的商已经属于 \(R[x]\)。因此 \(q_i\) 也为素元。

现在任一非零非单位都已分解为素元的乘积。比较两个不可约分解，第一个分解中的任一素因子必须整除第二个分解中的某个不可约因子，两者因而相伴；消去它们并重复，就得到因子个数相同且逐一相伴。这就是唯一分解性。有限多个不定元的结论，再用
\(R[x_1,\ldots,x_n]=\Bigl(R[x_1,\ldots,x_{n-1}]\Bigr)[x_n]\)
对 \(n\) 归纳。
\end{proof}

这个定理说明 \(\mathbb Z[x]\) 和 \(K[x,y]\) 都具有唯一分解性，却没有说明它们是主理想整环。唯一分解描述元素的因子分解；主理想条件要求一个理想中的全部元素由单个元素生成，是更强的要求。多元多项式环中理想的有限生成性，将由本章最后一节的 Hilbert 基定理保证。

% !TeX root = ../../Book1.tex
% 纲要标题：### 12.3 不可约性判据
% 纲要位置：计划纲要.md，第 456 行。

\section{不可约性判据}
\label{b1:sec:1203}

% 纲要内容（仅作写作计划，尚非教材正文）：
%
% * Eisenstein 判据及平移使用
% * 模素数约化法
% * 有理根检验的适用范围与失效例子
%

唯一分解定理保证因子分解存在而且唯一，却没有立即告诉我们某个给定多项式是否已经不可约。本节给出几种可以实际使用的判据。它们都只能在相应假设下推出结论；某一种检验没有奏效，并不等于多项式可约。

\subsection{Eisenstein 判据及平移}
\label{b1:subsec:1203:01}

设 \(R\) 为唯一分解整环，\(p\in R\) 为素元。若一个多项式模 \(p\) 后只剩最高次项，那么它的非平凡因子模 \(p\) 后也受到很强的限制。再观察常数项中 \(p\) 的幂次，就得到下面的判据。

\begin{theorem}[Eisenstein]\label{b1:thm:eisenstein}
设 \(f=a_nx^n+\cdots+a_1x+a_0\in R[x]\) 本原，\(n\ge1\)。若存在素元 \(p\)，满足
\[
p\nmid a_n,\qquad p\mid a_i\ (0\le i<n),\qquad p^2\nmid a_0,
\]
则 \(f\) 在 \(R[x]\) 及 \(\operatorname{Frac}(R)[x]\) 中均不可约。
\end{theorem}
\begin{proof}
由\cref{b1:thm:gauss-descent}，只需排除 \(R[x]\) 中两个正次数因子的分解。假设 \(f=gh\)，\(\deg g=r>0\)、\(\deg h=s>0\)，则 \(r+s=n\)。首项系数的积为 \(a_n\)，不被 \(p\) 整除，所以 \(g,h\) 的首项系数模 \(p\) 都非零；约化后次数仍为 \(r,s\)。

在整环 \(R/(p)\) 上，有 \(\bar g\bar h=\bar a_nx^n\)。设 \(\bar g,\bar h\) 中最低非零项的指数分别为 \(u,v\)。最低项系数的积非零，故乘积的最低指数为 \(u+v=n\)。另一方面 \(u\le r,v\le s,r+s=n\)，因此只能有 \(u=r,v=s\)。两者均为正数，所以 \(g,h\) 的常数项都被 \(p\) 整除。这使 \(a_0=g(0)h(0)\) 被 \(p^2\) 整除，与假设矛盾。
\end{proof}

这称为\term[eisenstein-panju]{Eisenstein 判据}[Eisenstein criterion]。例如 \(x^5-6x+3\) 对素数三满足该判据，故在 \(\mathbb Q[x]\) 中不可约。没有出现的系数就是零，当然也被三整除；常数项三不被九整除，是不能遗漏的条件。

\ProseParagraphBreak
有时原多项式的系数没有这种形式，平移后却有。对 \(a\in R\)，代入映射 \(T_a:R[x]\rightarrow R[x]\)、\(f(x)\mapsto f(x+a)\) 保持和、积与单位，而且 \(T_{-a}T_a=T_aT_{-a}=1\)。它是环自同构，因而保持单位与不可约性。这不是仅仅改变某个函数的图像，而是在多项式环内部作可逆代换。

\ProseParagraphBreak
以素数 \(p\) 为例，考虑
\[
\Phi_p(x)=1+x+\cdots+x^{p-1}.
\]
记号 \(\Phi_p\) 的一般推广将在后文介绍；这里仅指所列多项式。由几何级数恒等式，
\[
\Phi_p(x+1)=\frac{(x+1)^p-1}{x}
=\sum_{j=1}^p\binom pj x^{j-1}.
\]
右端说明商确为整数系数多项式。它的首项系数为一，常数项为 \(p\)；其余二项式系数都被 \(p\) 整除，因为 \(1\le j<p\) 时 \(j!\) 与 \((p-j)!\) 均不含素因子 \(p\)，而 \(p!\) 含有这个因子。故 \(\Phi_p(x+1)\) 对 \(p\) 满足 Eisenstein 判据，\(\Phi_p(x)\) 也在 \(\mathbb Q[x]\) 中不可约。

\subsection{模素数约化法}
\label{b1:subsec:1203:02}

把系数约化到有限域，可以把有理数域中的问题转成有限次的尝试。这个办法需要保住原多项式的次数。

\begin{proposition}\label{b1:prop:reduction-irreducible}
设 \(f\in\mathbb Z[x]\) 本原，次数为正，\(p\) 是不整除其首项系数的素数。若逐系数约化所得的 \(\bar f\in\mathbb F_p[x]\) 不可约，则 \(f\) 在 \(\mathbb Q[x]\) 及 \(\mathbb Z[x]\) 中都不可约。
\end{proposition}
\begin{proof}
若 \(f\) 在 \(\mathbb Q[x]\) 中可约，\cref{b1:thm:gauss-descent}给出 \(f=gh\)，其中 \(g,h\in\mathbb Z[x]\) 的次数均为正。由于首项系数的积不被 \(p\) 整除，每个因子的首项系数也不被 \(p\) 整除，所以 \(\deg\bar g=\deg g>0\)、\(\deg\bar h=\deg h>0\)。约化后 \(\bar f=\bar g\bar h\) 就是非平凡分解，与假设矛盾。再由同一个 Gauss 分解定理，得到 \(\mathbb Z[x]\) 中的结论。
\end{proof}

上述方法称为\term[mo-sushu-yuehua-fa]{模素数约化法}[reduction modulo a prime]。例如 \(f=x^3+x+1\) 模二后，在 \(\bar0,\bar1\) 处的值都为 \(\bar1\)，所以没有根。域上的二次或三次多项式若可约，两个正次数因子的次数之和至多为三，其中必有一次因子；由\cref{b1:prop:factor-root}，这等价于存在域内的根。因此 \(\bar f\) 不可约，所给整数多项式在 \(\mathbb Q[x]\) 中不可约。

\ProseParagraphBreak
若约化后可约，只能说这次检验没有得出结论。例如 \(x^2+1\) 在 \(\mathbb Q[x]\) 中不可约，因为有理数的平方不可能等于负一；但它模二后等于 \((x+1)^2\)。首项系数条件也不能删去：
\[
2x^2+3x+1=(2x+1)(x+1)
\]
模二后却成为一次不可约多项式 \(x+1\)。失去的因子被约化成了常数，原来的次数分配便无法恢复。

\subsection{有理根检验的适用范围}
\label{b1:subsec:1203:03}

\begin{proposition}[有理根检验]\label{b1:prop:rational-root}
设 \(f=a_nx^n+\cdots+a_0\in\mathbb Z[x]\)，\(a_n\ne0\)。若非零有理数 \(r/s\) 是 \(f\) 的根，其中 \(r,s\in\mathbb Z\)、\(s>0\)、\(\gcd(r,s)=1\)，则 \(r\mid a_0\)、\(s\mid a_n\)。特别地，首一整系数多项式的每个有理根都是整数。
\end{proposition}
\begin{proof}
把 \(f(r/s)=0\) 乘以 \(s^n\)，得到
\[
a_nr^n+a_{n-1}r^{n-1}s+\cdots+a_1rs^{n-1}+a_0s^n=0.
\]
除最后一项以外，每项都被 \(r\) 整除，因此 \(r\mid a_0s^n\)。由于 \(r\) 与 \(s^n\) 互素，由整数的素因子分解可消去 \(s^n\)，得到 \(r\mid a_0\)。同样，除第一项以外，每项都被 \(s\) 整除，故 \(s\mid a_nr^n\)；互素性给出 \(s\mid a_n\)。若 \(a_n=1\)，正分母只能为一。
\end{proof}

这称为\term[youli-gen-jianyan]{有理根检验}[rational root test]。零是否为根，直接查看 \(a_0\) 是否为零即可。当 \(a_0\ne0\) 时，定理只留下有限个可能的分子和分母，可以逐一代入。例如 \(x^3-2\) 的候选有理根只有 \(\pm1,\pm2\)，代入均不为零；次数为三，故它在 \(\mathbb Q[x]\) 中不可约。这里当然也可以直接使用素数二的 Eisenstein 判据。

\ProseParagraphBreak
次数达到四以后，没有有理根只能排除一次因子，不能排除两个二次因子的乘积。例如
\[
x^4+x^2+1=(x^2+x+1)(x^2-x+1).
\]
左边对任何实数都严格为正，因而没有有理根，却在 \(\mathbb Q[x]\) 中可约。选择判据时，先问它排除了哪些可能的因子次数，往往比机械地尝试更多代入值更有效。

% !TeX root = ../../Book1.tex
% 纲要标题：### 12.4 多元多项式与有限性
% 纲要位置：计划纲要.md，第 462 行。

\section{多元多项式与有限性}
\label{b1:sec:1204}

% TODO: 按以下纲要要点撰写本节。

% 纲要内容（仅作写作计划，尚非教材正文）：
%
% * 单项式、多项式次数及消去的动机
% * Noether 环与升链条件
% * Hilbert 基定理；第三卷只复用证明并发展模论版本
%
% ---
%

待撰写。


\begin{chaptersummary}
因式定理把根转化为一次因子，形式导数把重根转化为最大公因式问题；插值则用次数界保证有限个取值条件的唯一性。Gauss 引理从素元约化出发，证明容度的乘法性质，使分式域上的非平凡分解下降到本原多项式的系数环，并进一步建立多项式环的唯一分解性。不可约性检验始终依赖所取的基域和多项式次数，尤其不能把“没有有理根”当作高次多项式不可约的充分条件。

\ProseParagraphBreak
交换变量而保持不变的多项式，可以唯一写成初等对称多项式的多项式；首项消去在有界全次数的有限单项式集合中终止，这一证明适用于任意交换系数环。根与系数的关系由此成为处理对称表达式的具体工具。Hilbert 基定理处理的是另一类有限性：有限多个不定元不会破坏 Noether 性，而其证明并未要求底环有唯一分解。本章已经构造分式域，并用它比较多项式分解；下一章将研究这一构造的映射性质，再讨论只允许指定元素作分母的局部化。
\end{chaptersummary}

\BookExercises

\begin{exercise}\label{b1:ex:12-polynomial-functions}
设 \(p\) 为素数。证明 \(x^p-x\) 在 \(\mathbb F_p\) 的每个元素处取值为零，但不是零多项式。进而证明任意函数 \(\varphi:\mathbb F_p\rightarrow\mathbb F_p\) 都由唯一一个次数至多为 \(p-1\) 的多项式给出，并确定所有表示同一函数的多项式之间的关系。
\end{exercise}
\begin{solution}
若 \(a=0\)，等式 \(a^p-a=0\) 直接成立；若 \(a\ne0\)，它属于 \(p-1\) 阶乘法群 \(\mathbb F_p^\times\)。\cref{b1:thm:lagrange}说明 \(a\) 的阶整除 \(p-1\)，所以 \(a^{p-1}=1\)，仍得 \(a^p=a\)。多项式 \(x^p-x\) 的首项系数为一，故不是零多项式。

将 \(\mathbb F_p\) 中全部 \(p\) 个元素作为插值点，给定值取为 \(\varphi\) 的值，\cref{b1:thm:lagrange-interpolation}就给出次数至多为 \(p-1\) 的唯一代表。若 \(f,g\) 给出同一个函数，对 \(f-g\) 除以 \(x^p-x\)，得到 \(f-g=q(x^p-x)+r\)，其中 \(r=0\) 或 \(\deg r<p\)。前两项在每个点都为零，因此 \(r\) 在 \(p\) 个点都为零。\cref{b1:prop:factor-root}迫使 \(r=0\)，所以 \(x^p-x\mid f-g\)。反过来，该整除关系显然保证两者的全部取值相同。故它们恰好相差理想 \((x^p-x)\) 中的元素。
\end{solution}

\begin{exercise}\label{b1:ex:12-repeated-roots}
在 \(\mathbb Q[x]\) 中求
\[
f=x^4-2x^3+2x^2-2x+1
\]
与 \(f'\) 的首一最大公因式，并给出它在 \(\mathbb C\) 中各根的重数。再将系数模二约化，比较重根情况。
\end{exercise}
\begin{solution}
直接展开可验
\[
f=(x-1)^2(x^2+1),\qquad
f'=2(x-1)(2x^2-x+1).
\]
其中 \(2x^2-x+1\) 在 \(x=1\) 处的值为二，所以不能再提出 \(x-1\)。它与 \(x^2+1\) 的最大公因式也为一：两者相减适当倍数后得到 \(-x-1\)，而 \(x^2+1\) 在负一处的值为二，故不能含有这个一次因子。因此首一最大公因式为 \(x-1\)。在 \(\mathbb C[x]\) 中，\(f=(x-1)^2(x-i)(x+i)\)，三个根互异，所以一的重数为二，\(i,-i\) 的重数均为一。

模二后 \(\bar f=x^4+1=(x+1)^4\)，形式导数为零。因此在任何包含 \(\mathbb F_2\) 的域中，它只有根一，重数为四。约化不仅可能使不同的因子合并，也可能使导数中的系数全部消失，不能照搬特征零下的重数计算。
\end{solution}

\begin{exercise}\label{b1:ex:12-interpolation}
求次数至多为 \(2\)，且满足 \(f(0)=1,f(1)=0,f(2)=3\) 的唯一有理系数多项式。再确定所有满足这些取值条件的 \(\mathbb Q[x]\) 中的多项式。
\end{exercise}
\begin{solution}
在\eqref{b1:eq:lagrange-interpolation}中代入这三个点和值，得到
\[
f_0(x)=\frac{(x-1)(x-2)}2+\frac{3x(x-1)}2
=2x^2-3x+1.
\]
若 \(f\) 满足相同的条件，\(f-f_0\) 在零、一、二处均为零。由\cref{b1:prop:factor-root}，先提出 \(x\)，所得商在一、二处仍为零，因为在这两点 \(x\) 的值非零；再依次提出 \(x-1,x-2\)，得到
\[
f=f_0+x(x-1)(x-2)h,\qquad h\in\mathbb Q[x].
\]
反向代入即可验证每个这样的多项式都符合条件。若还要求次数至多为二，则只能 \(h=0\)，这也直接验证了唯一性。
\end{solution}

\begin{exercise}\label{b1:ex:12-content-factorization}
沿用正文对 \(6x^2+9x+3\) 的容度分解，分别在 \(\mathbb Z[x]\) 与 \(\mathbb Q[x]\) 中验证所列因子确为不可约因子，并比较两个环中的单位部分与常数因子。
\end{exercise}
\begin{solution}
容度为三，本原部分为 \(2x^2+3x+1\)，而
\[
6x^2+9x+3=3(2x+1)(x+1).
\]
在 \(\mathbb Z[x]\) 中，三为素元：逐系数模三约化的核为 \((3)\)，像为整环 \(\mathbb F_3[x]\)，所以若三整除某个乘积，约化后必有一个因子为零，即三整除这个因子的全部系数。两个一次多项式均本原，在 \(\mathbb Q[x]\) 中因次数为一而不可约，再由\cref{b1:thm:gauss-descent}，在 \(\mathbb Z[x]\) 中也不可约。故所列三因子就是不可约分解，单位只能是 \(\pm1\)。

在 \(\mathbb Q[x]\) 中，所有非零有理常数都是单位，可以写成
\(6(x+\tfrac12)(x+1)\)。这里只有两个正次数不可约因子，常数六属于单位部分。唯一分解的比较必须在同一个环中进行；扩大系数环以后，原来的某些非单位会变为单位。
\end{solution}

\begin{exercise}\label{b1:ex:12-monic-descent}
设首一多项式 \(f\in\mathbb Z[x]\) 在 \(\mathbb Q[x]\) 中写成 \(f=gh\)，其中 \(g,h\) 也首一。证明 \(g,h\in\mathbb Z[x]\)。再判断以下两种变式能否保证因子仍有整数系数：只要求 \(f\) 本原，并允许写成 \(f=cgh\)，其中 \(c\in\mathbb Q^\times\)、\(g,h\) 首一；或者保持 \(f=gh\) 与 \(f\) 首一，却取消两个因子的首一要求。
\end{exercise}
\begin{solution}
首一的 \(f\) 本原。由\cref{b1:lem:primitive-scaling}，可以写
\[
g=\alpha g_0,\qquad h=\beta h_0,
\]
其中 \(g_0,h_0\in\mathbb Z[x]\) 本原，\(\alpha,\beta\in\mathbb Q^\times\)。\cref{b1:lem:gauss-primitive}说明 \(g_0h_0\) 本原；而 \(f=\alpha\beta g_0h_0\) 也本原，所以\cref{b1:lem:primitive-scaling}给出 \(\alpha\beta=\pm1\)。设 \(g_0,h_0\) 的首项系数为非零整数 \(a,b\)。比较 \(f\) 的首项系数，得 \(\alpha\beta ab=1\)，因而 \(ab=\pm1\)，所以 \(a,b\) 都等于 \(\pm1\)。再分别比较首一的 \(g,h\) 的首项系数，得 \(\alpha=1/a\)、\(\beta=1/b\)，两者也是整数单位。因此 \(g,h\) 的全部系数都是整数。

第一种变式不能保证结论。例如
\[
2x^2+3x+1=2\Bigl(x+\frac12\Bigr)(x+1).
\]
如果保持 \(f\) 首一，却取消两个因子的首一要求，结论同样失败：\(x^2=(2x)(x/2)\)。因子的非零常数倍可以互相转移，首一条件恰好排除了这种任意缩放。
\end{solution}

\begin{exercise}\label{b1:ex:12-shift-eisenstein}
证明 \(x^4+1\) 在 \(\mathbb Q[x]\) 中不可约，并说明对原多项式直接使用 Eisenstein 判据为何行不通。再证明它对每个素数 \(p\) 的约化在 \(\mathbb F_p[x]\) 中都可约。
\end{exercise}
\begin{solution}
对 \(f=x^4+1\)，常数项为一，不被任何素数整除，所以原式不符合 Eisenstein 判据。作可逆平移得
\[
f(x+1)=x^4+4x^3+6x^2+4x+2.
\]
首项系数为一，其余系数都被二整除，常数项二不被四整除。由\cref{b1:thm:eisenstein}，\(f(x+1)\) 在 \(\mathbb Q[x]\) 中不可约。若 \(f\) 有非平凡分解，逐因子代入 \(x+1\) 会保持各因子的正次数，给出 \(f(x+1)\) 的非平凡分解，矛盾。因此 \(f\) 不可约。

模二时 \(x^4+1=(x+1)^4\)，所以约化式可约。设 \(p\) 为奇素数。非零平方在 \(\mathbb F_p^\times\) 中组成指数为二的子群：每个非零平方恰有两个平方根 \(a,-a\)。若 \(-1\) 是平方，取 \(a^2=-1\)，则
\[
x^4+1=(x^2-a)(x^2+a).
\]
若 \(-1\) 不是平方，则乘以这个非平方元会交换平方子群及其非平方陪集，故 \(2\) 与 \(-2\) 中恰有一个是平方。若 \(a^2=2\)，则
\[
x^4+1=(x^2+ax+1)(x^2-ax+1);
\]
若 \(a^2=-2\)，则
\[
x^4+1=(x^2+ax-1)(x^2-ax-1).
\]
每种情形都给出了两个二次因子。因此此例虽在 \(\mathbb Q[x]\) 中不可约，却找不到任何不可约的模素数约化；逐个尝试素数的办法在此不会成功。
\end{solution}

\begin{exercise}\label{b1:ex:12-quartic-mod-two}
用模二约化证明 \(x^4+x+1\) 在 \(\mathbb Q[x]\) 中不可约。论证中除去一次因子的可能性后，还需要检查什么？
\end{exercise}
\begin{solution}
约化多项式在 \(\bar0,\bar1\) 处均取 \(\bar1\)，所以没有一次因子。若它可约，按不可约因子的次数分拆四，只能是两个不可约二次因子的乘积。\(\mathbb F_2[x]\) 中的首一不可约二次多项式只有 \(x^2+x+1\)：常数项必须为一，以免零为根；对 \(x^2+ax+1\) 代入一得到 \(a\)，所以还必须 \(a=1\)。但是
\[
(x^2+x+1)^2=x^4+x^2+1\ne x^4+x+1
\]
于 \(\mathbb F_2[x]\)。因此二次乘二次的可能性也排除了，约化式不可约。原多项式本原且首项系数为一，\cref{b1:prop:reduction-irreducible}于是给出 \(\mathbb Q[x]\) 中的不可约性。
\end{solution}

\begin{exercise}\label{b1:ex:12-rational-root-limits}
用有理根检验分解 \(6x^3+5x^2-12x+4\) 为 \(\mathbb Q[x]\) 中的不可约因子。再证明 \(x^4+4\) 没有有理根，却在 \(\mathbb Q[x]\) 中可约，并写出其不可约分解。
\end{exercise}
\begin{solution}
由\cref{b1:prop:rational-root}，第一个多项式的非零有理根约成 \(r/s\)、\(s>0\) 后，必须满足 \(r\mid4\)、\(s\mid6\)。从这些候选中代入 \(1/2\)，得到零。除以 \(2x-1\)，有
\[
6x^3+5x^2-12x+4=(2x-1)(3x^2+4x-4).
\]
对剩下的二次多项式，检验 \(2/3\) 得到零，继续相除得
\[
6x^3+5x^2-12x+4=(2x-1)(3x-2)(x+2).
\]
三个因子在 \(\mathbb Q[x]\) 中都是一次多项式，所以均不可约；它们给出了全部三个根 \(1/2,2/3,-2\)。

对任何有理数 \(a\)，都有 \(a^4+4>0\)，所以第二个多项式没有有理根。但将其写成平方差，可得
\[
x^4+4=(x^2+2)^2-(2x)^2
=(x^2-2x+2)(x^2+2x+2).
\]
两个二次因子分别为 \((x-1)^2+1\)、\((x+1)^2+1\)，因此都没有有理根。二次多项式如果可约就必须有一次因子，而\cref{b1:prop:factor-root}会给出一个有理根；所以这两个二次因子均不可约。原式没有有理根，只排除了一次因子，并没有排除这里的二次因子。
\end{solution}

\begin{exercise}\label{b1:ex:12-primitive-linear-polynomials}
设 \(k\) 为任意域。分别判断 \(yx+y+1\) 和 \(yx+y\) 在 \(k[x,y]\) 中是否不可约。说明“关于 \(x\) 的次数为一”为什么不足以直接作出判断。再设 \(n\ge2\)，把 \(y\) 看作系数环 \(k[y]\) 中的元素，证明 \(x^n-y\) 在 \(k[x,y]\) 和 \(\operatorname{Frac}(k[y])[x]\) 中都不可约。若 \(\operatorname{char}k=p>0\) 且 \(n=p\)，举出一个包含 \(\operatorname{Frac}(k[y])\) 的域，使这个不可约多项式在其中有重根。
\end{exercise}
\begin{solution}
把两个多项式看作 \(k[y][x]\) 中的元素。\(k[y]\) 是唯一分解整环，因为它是域上一元多项式环，适用\cref{b1:thm:polynomial-ufd}。第一个多项式的两个系数 \(y,y+1\) 互素：任何公因子都整除它们之差一。因此 \(yx+y+1\) 本原。进入 \(\operatorname{Frac}\Bigl(k[y]\Bigr)[x]\) 后，它是一次多项式，故不可约；再由\cref{b1:thm:gauss-descent}，它在 \(k[y][x]=k[x,y]\) 中也不可约。

第二个多项式的容度可取 \(y\)，并且 \(yx+y=y(x+1)\)。这里两个因子都是非单位：若它们有多项式逆元，非零乘积的全次数相加将迫使它们的次数为零，矛盾。所以它可约。\(y\) 虽然关于 \(x\) 是常数，在系数环 \(k[y]\) 中却不是单位；这正是不能仅凭关于 \(x\) 的次数来套用域上结论的原因。

对第三个多项式，\(k[y]\) 是唯一分解整环，\(y\) 是其中的素元。把 \(x^n-y\) 看作 \(k[y][x]\) 中关于 \(x\) 的多项式：首项系数为一，因此它本原；中间各项系数为零，常数项 \(-y\) 被 \(y\) 整除但不被 \(y^2\) 整除。由\cref{b1:thm:eisenstein}，它在 \(\operatorname{Frac}(k[y])[x]\) 中不可约。再由\cref{b1:thm:gauss-descent}，它在 \(k[x,y]\) 中也不可约。

若 \(n=p=\operatorname{char}k\)，记 \(F=\operatorname{Frac}(k[y])\)，并取
\[
L=F[t]/(t^p-y),\qquad \alpha=t+(t^p-y)\in L.
\]
已证 \(t^p-y\) 在 \(F[t]\) 中不可约，故 \(L\) 是域，且 \(\alpha^p=y\)。在 \(L[x]\) 中，特征为 \(p\)，于是
\[
x^p-y=x^p-\alpha^p=(x-\alpha)^p.
\]
所以在 \(F[x]\) 中不可约并不妨碍它在扩大系数域后出现重根。
\end{solution}

\begin{exercise}\label{b1:ex:12-cyclic-polynomial-invariants}
设 \(k\) 为域，三阶循环群 \(C_3=\langle\sigma\rangle\) 通过 \(\sigma\cdot x=y\)、\(\sigma\cdot y=z\)、\(\sigma\cdot z=x\) 作用在 \(k[x,y,z]\) 上。列出全次数至多为二的全部不变多项式，并说明它们是否一定对称。再给出一个三次齐次不变多项式，它不是对称多项式。你的结论是否需要排除特征三？
\end{exercise}
\begin{solution}
次数至多为二的单项式在这个作用下分为四个轨道：
\[
\{1\},\qquad \{x,y,z\},\qquad
\{x^2,y^2,z^2\},\qquad \{xy,yz,zx\}.
\]
比较单项式系数，不变性恰好要求每个轨道中的系数相等。因此全部所求多项式唯一写成
\[
a+b(x+y+z)+c(x^2+y^2+z^2)+d(xy+yz+zx),
\qquad a,b,c,d\in k.
\]
每个括号内的表达式都在任意变量置换下不变，所以次数至多为二时，循环不变性确实蕴含对称性。

三次时取 \(f=x^2y+y^2z+z^2x\)。循环置换重排这三项，故 \(f\) 不变；交换 \(x,y\) 得到 \(xy^2+yz^2+zx^2\)，其中 \(x^2y\) 的系数为零，原式中则为一，所以两式不同。这给出所需反例。全部论证只比较系数，从未除以群阶，故在特征三的域上同样成立。
\end{solution}

\begin{exercise}\label{b1:ex:12-symmetric-fourth-power}
设 \(R\) 为交换环，在 \(R[x,y,z]\) 中记 \(e_1=x+y+z\)、\(e_2=xy+xz+yz\)、\(e_3=xyz\)。把 \(x^4+y^4+z^4\) 写成 \(e_1,e_2,e_3\) 的多项式，不作任何除法。说明结果在特征二时怎样简化，以及为什么这种简化不违背对称表达式的唯一性。
\end{exercise}
\begin{solution}
记 \(p_j=x^j+y^j+z^j\)。展开平方以及\eqref{b1:eq:power-sum-three}分别给出
\[
p_1=e_1,\qquad p_2=e_1^2-2e_2,\qquad
p_3=e_1^3-3e_1e_2+3e_3.
\]
由\eqref{b1:eq:elementary-symmetric-product}，代入 \(T=x\)、\(y\)、\(z\)，每个变量 \(u\) 都满足 \(u^3-e_1u^2+e_2u-e_3=0\)。分别乘以 \(u\) 后相加，得到
\[
\begin{aligned}
p_4&=e_1p_3-e_2p_2+e_3p_1\\
&=e_1^4-4e_1^2e_2+2e_2^2+4e_1e_3.
\end{aligned}
\]
全程只有环中的加法和乘法，所以对所给任意交换环成立。在特征二时，二和四的系数都为零，因而 \(p_4=e_1^4\)。唯一性是在固定系数环 \(R\) 中比较多项式；此时
\(t_1^4-4t_1^2t_2+2t_2^2+4t_1t_3\) 与 \(t_1^4\) 本来就是 \(R[t_1,t_2,t_3]\) 中的同一个多项式，并没有给出两个不同表达式。
\end{solution}

\begin{exercise}\label{b1:ex:12-invariants-fixed-variable}
设 \(k\) 为域，\(S_2\) 在 \(k[x,y,z]\) 上交换 \(x,y\) 而固定 \(z\)。证明每个不变多项式唯一写成 \(F(x+y,xy,z)\)，其中 \(F\in k[s,p,t]\)。进一步，令 \(S_2\times S_2\) 分别交换 \(x_1,x_2\) 和 \(y_1,y_2\)，求它在 \(k[x_1,x_2,y_1,y_2]\) 上的不变环，并证明所取生成元之间没有非零的多项式关系。
\end{exercise}
\begin{solution}
第一问把 \(z\) 放进系数环 \(R=k[z]\)。一个多项式在交换 \(x,y\) 后不变，恰好说明它是 \(R[x,y]\) 中的对称多项式。\cref{b1:thm:fundamental-symmetric-polynomials}给出唯一的 \(H\in R[s,p]\)，满足 \(f=H(x+y,xy)\)。将 \(H\) 的各系数唯一展开为 \(z\) 的多项式，就得到所需的 \(F\in k[s,p,t]\)；这两个唯一性合在一起也证明 \(F\) 唯一。

第二问记
\[
s_x=x_1+x_2,\quad p_x=x_1x_2,\qquad
s_y=y_1+y_2,\quad p_y=y_1y_2.
\]
先要求交换 \(x_1,x_2\) 后不变，应用同一定理，得到唯一表达式
\[
f=\sum_{a,b}c_{ab}(y_1,y_2)s_x^ap_x^b,
\qquad c_{ab}\in k[y_1,y_2],
\]
其中只有有限多个系数非零。再交换 \(y_1,y_2\)。由于关于 \(s_x,p_x\) 的表达式唯一，第二个交换也固定 \(f\)，当且仅当每个 \(c_{ab}\) 都对称。再次使用同一定理，每个系数唯一写成 \(s_y,p_y\) 的多项式。故不变环为
\[
k[x_1,x_2,y_1,y_2]^{S_2\times S_2}
=k[s_x,p_x,s_y,p_y].
\]
若这四个生成元满足一个非零多项式关系，先按前两个变量的幂收集，第一轮唯一性迫使每个关于 \(s_y,p_y\) 的系数为零；第二轮唯一性又迫使这些系数多项式的所有系数为零，与关系非零矛盾。因而它们之间没有非零多项式关系。
\end{solution}

\begin{exercise}\label{b1:ex:12-squared-root-elimination}
设 \(L\) 为包含 \(\mathbb Q\) 的域，并且在 \(L[X]\) 中有分解
\[
X^3-2X^2+3X-4=(X-a)(X-b)(X-c).
\]
不逐个求出 \(a,b,c\)，确定首一多项式 \((T-a^2)(T-b^2)(T-c^2)\)，并证明其系数属于 \(\mathbb Q\)。
\end{exercise}
\begin{solution}
记 \(e_1=a+b+c\)、\(e_2=ab+ac+bc\)、\(e_3=abc\)。比较给定分解的系数得 \(e_1=2\)、\(e_2=3\)、\(e_3=4\)。展开可得
\[
\begin{aligned}
a^2+b^2+c^2&=e_1^2-2e_2=-2,\\
a^2b^2+a^2c^2+b^2c^2&=e_2^2-2e_1e_3=-7,\\
a^2b^2c^2&=e_3^2=16.
\end{aligned}
\]
其中第二式的依据是 \((ab+ac+bc)^2\) 的交叉项之和为 \(2abc(a+b+c)\)。再次使用根与系数关系，得到
\[
(T-a^2)(T-b^2)(T-c^2)=T^3+2T^2-7T-16.
\]
这些系数都是 \(e_1,e_2,e_3\) 的整数系数多项式，而三个 \(e_i\) 都是有理数，故所求多项式属于 \(\mathbb Q[T]\)。这里不需要知道各个根属于哪个具体扩域，也不需要判断平方后的三个根是否互异。
\end{solution}

\begin{exercise}\label{b1:ex:12-noetherian-not-pid}
证明 \(\mathbb Q[x,y]\) 是 Noether 唯一分解整环，却不是主理想整环。再证明 \(\mathbb Q[x,y]/(xy)\) 是 Noether 环，但不是整环。
\end{exercise}
\begin{solution}
域 \(\mathbb Q\) 是 Noether 唯一分解整环，故\cref{b1:thm:polynomial-ufd}与\cref{b1:thm:hilbert-basis}分别给出 \(\mathbb Q[x,y]\) 的唯一分解性和 Noether 性。理想 \((x,y)\) 为真理想，因为其中每个多项式在 \((0,0)\) 处为零。如果 \((x,y)=(g)\)，则 \(g\mid x\) 且 \(g\mid y\)。用全次数的乘法公式，由 \(g\mid x\) 知 \(g\) 是非零常数或与 \(x\) 相伴；后者不可能整除 \(y\)，因为代入 \(x=0\) 会给出 \(y=0\)。因此 \(g\) 为非零有理常数，是单位，与理想为真矛盾。所以该环不是主理想整环。

由\cref{b1:prop:noetherian-quotient}，商环 \(\mathbb Q[x,y]/(xy)\) 为 Noether 环。\(x,y\) 的剩余类都非零：若 \(x\in(xy)\)，代入 \(y=0\) 就得 \(x=0\)，矛盾；对 \(y\) 代入 \(x=0\) 同理。但两者的积为零，故该商环不是整环。
\end{solution}

\begin{exercise}\label{b1:ex:12-elimination}
在 \(\mathbb Q[x,y,z]\) 中令 \(I=(y-x^2,z-xy,z-1)\)。求 \(I\cap\mathbb Q[x]\)，并解释它如何限制原方程组的第一坐标。
\end{exercise}
\begin{solution}
直接计算
\[
x^3-1=(z-1)-(z-xy)-x(y-x^2),
\]
说明 \(x^3-1\in I\)。反过来，
\[
z-xy=(z-1)-(x^3-1)-x(y-x^2),
\]
所以 \(I=(y-x^2,z-1,x^3-1)\)。若 \(g(x)\in I\)，将它表示成这三个生成元的多项式系数组合，再代入 \(y=x^2,z=1\)，得到 \(g(x)=C(x,x^2,1)(x^3-1)\)。故 \(g\) 是 \(x^3-1\) 的倍数。结合此前已经证明的包含，得到
\[
I\cap\mathbb Q[x]=(x^3-1).
\]
原方程组的任一解必满足 \(x^3=1\)。本题中反向提升也成立：在任意包含 \(\mathbb Q\) 的域中，若 \(a^3=1\)，则 \((a,a^2,1)\) 同时满足三个原方程，且后两个坐标由第一个坐标唯一确定。这一充分性来自本题明确的代入，不是消去理想在一般情况下自动保证的结论。
\end{solution}

\begin{exercise}\label{b1:ex:12-noetherian-surjection}
证明 Noether 交换环 \(R\) 的每个满环自同态 \(\varphi:R\rightarrow R\) 都是同构。用无穷多个不定元的多项式环，说明 Noether 条件不能删去。
\end{exercise}
\begin{solution}
记 \(\varphi^n\) 为 \(\varphi\) 的 \(n\) 次复合。各次复合的核组成理想升链
\[
\ker\varphi\subseteq\ker\varphi^2\subseteq\cdots.
\]
因为若 \(\varphi^n(a)=0\)，再作用一次 \(\varphi\) 仍为零，所以这些包含成立。\cref{b1:prop:noetherian-acc}给出某个 \(N\ge1\)，使 \(\ker\varphi^N=\ker\varphi^{N+1}\)。任取 \(a\in\ker\varphi\)，由于 \(\varphi\) 满射，\(\varphi^N\) 也满射，可以取 \(b\in R\) 使 \(\varphi^N(b)=a\)。于是
\[
\varphi^{N+1}(b)=\varphi(a)=0,
\]
故 \(b\in\ker\varphi^{N+1}=\ker\varphi^N\)，从而 \(a=0\)。这证明 \(\varphi\) 单射，结合满射性即为同构。

去掉 Noether 条件，取 \(R=k[x_1,x_2,\ldots]\)，令 \(\varphi\) 固定 \(k\)，把 \(x_1\) 送到零，把 \(x_{i+1}\) 送到 \(x_i\)。每个多项式只含有限多个变量，这些代入确定保持加法、乘法和单位的映射。任意多项式的原像可以把其中每个 \(x_i\) 换成 \(x_{i+1}\) 得到，所以映射满射；非零元素 \(x_1\) 却在其核中，故不单射。相应核升链确实不稳定：\(x_{n+1}\) 在 \(\ker\varphi^{n+1}\) 中，却不在 \(\ker\varphi^n\) 中，因为 \(\varphi^n(x_{n+1})=x_1\ne0\)。
\end{solution}

\begin{exercise}\label{b1:ex:12-nonnoetherian-subring}
设 \(k\) 为域，在 \(k[x,y]\) 中取子环
\[
A=k[x,xy,xy^2,xy^3,\ldots].
\]
证明 \(A=k+xk[x,y]\)，并证明 \(A\) 不是 Noether 环，尽管 \(k[x,y]\) 是。这里的 \(k+xk[x,y]\) 表示所有 \(c+xh(x,y)\)、\(c\in k\)、\(h\in k[x,y]\) 组成的集合。
\end{exercise}
\begin{solution}
所列每个生成元都属于 \(k+xk[x,y]\)，而后者在加法、取负和乘法下封闭，故 \(A\subseteq k+xk[x,y]\)。反过来，任意 \(x^iy^j\)、\(i\ge1,j\ge0\)，都可写成 \(x^{i-1}(xy^j)\)，属于 \(A\)。每个 \(xh(x,y)\) 是这种单项式的有限线性组合，所以另一个包含也成立。

在 \(A\) 中令 \(I_n=(x,xy,\ldots,xy^n)\)。这些理想构成升链。如果 \(xy^{n+1}\in I_n\)，则存在 \(a_0,\ldots,a_n\in A\)，使
\[
xy^{n+1}=\sum_{j=0}^n a_jxy^j.
\]
根据刚证的环的描述，写 \(a_j=c_j+xh_j(x,y)\)，其中 \(c_j\in k\)。比较两边作为 \(k[y][x]\) 中多项式的 \(x\) 的一次项系数，便得到
\[
y^{n+1}=\sum_{j=0}^n c_jy^j,
\]
这与多项式系数的唯一性矛盾。所以每个 \(I_n\subsetneq I_{n+1}\)，由\cref{b1:prop:noetherian-acc}，\(A\) 不是 Noether 环。另一方面，域 \(k\) 是 Noether 环，\cref{b1:thm:hilbert-basis}保证 \(k[x,y]\) 是 Noether 环。因此 Noether 性并不任意传给子环。
\end{solution}

\begin{exercise}\label{b1:ex:12-closed-field-evaluation}
设 \(k\) 为代数闭域，\(f\in k[x]\) 首一且次数为 \(n\ge1\)，其不同的根为 \(a_1,\ldots,a_r\)。证明取值公式给出满环同态
\[
E:k[x]/(f)\longrightarrow k^r,\qquad
g+(f)\longmapsto\Bigl(g(a_1),\ldots,g(a_r)\Bigr),
\]
并证明 \(\ker E\) 恰为商环中全部幂零元组成的集合。最后证明 \(E\) 是同构，当且仅当 \(f\) 没有重根。
\end{exercise}
\begin{solution}
由\cref{b1:prop:algebraically-closed-splitting}，存在正整数 \(m_1,\ldots,m_r\)，使
\[
f=\prod_{i=1}^r(x-a_i)^{m_i},\qquad
\sum_{i=1}^r m_i=n.
\]
如果 \(g-g'\in(f)\)，那么在每个 \(a_i\) 处两者取值相同，所以 \(E\) 良定义。求值保持加法、乘法和单位，故 \(E\) 是环同态。\cref{b1:thm:lagrange-interpolation}说明任意指定的 \(r\) 个值都可由一个次数至多为 \(r-1\) 的多项式取得，因此 \(E\) 满射。

若 \(z=g+(f)\) 幂零，存在 \(N\ge1\) 使 \(g^N\in(f)\)。代入每个 \(a_i\) 得 \(g(a_i)^N=0\)；域中没有非零幂零元，所以 \(g(a_i)=0\)，即 \(z\in\ker E\)。反过来，若这些取值全为零，\cref{b1:prop:factor-root}保证每个 \(x-a_i\) 都整除 \(g\)。这些一次因子两两不相伴，而\cref{b1:thm:polynomial-ufd}保证 \(k[x]\) 的唯一分解性，所以其乘积
\(h=\prod_{i=1}^r(x-a_i)\) 因而整除 \(g\)。取 \(N=\max_i m_i\)，则 \(f\mid h^N\mid g^N\)，所以 \(z^N=0\)。这就刻画了核。

如果 \(f\) 没有重根，所有 \(m_i=1\)，于是 \(f=h\)。刚才的论证说明 \(E\Bigl(g+(f)\Bigr)=0\) 就蕴含 \(f\mid g\)，故 \(E\) 单射，从而为同构。如果至少一个 \(m_i>1\)，则 \(r<n\)，多项式 \(h\) 的次数为 \(r\)，小于 \(f\) 的次数，所以 \(h\notin(f)\)；但 \(h\) 在各 \(a_i\) 处均为零。因此 \(h+(f)\) 是核中的非零元素，\(E\) 不是同构。
\end{solution}
